\begin{table}[tbp]\centering\small
\caption{Unemployment Response to a Tightening of the Market Index}\label{tab:var}
\resizebox{\ifdim\width>\textwidth\textwidth\else\width\fi}{!}{%
\begin{tabular}{lrrrrrrr}
\toprule
 & Lags & Granger $F$ & $p$ & $\Phi(6)$ & $\Phi(12)$ [90\%] & $\Phi(24)$ & Months significant \\
\midrule
Nominal & 12 & 3.74 & $<$0.0001 & 0.29 & 0.57 [0.29, 0.77] & 0.66 & 4--36 \\
Real & 12 & 3.67 & $<$0.0001 & 0.34 & 0.61 [0.32, 0.82] & 0.68 & 3--36 \\
Relative to $r^*$ (baseline) & 12 & 5.30 & $<$0.0001 & 0.40 & 0.65 [0.39, 0.83] & 0.64 & 1--36 \\
Relative to $r^*$ (Kiley) & 12 & 4.63 & $<$0.0001 & 0.37 & 0.61 [0.33, 0.80] & 0.63 & 3--36 \\
Relative to $r^*$, 4 lags & 4 & 9.98 & $<$0.0001 & 0.32 & 0.40 [0.21, 0.58] & 0.38 & 1--36 \\
\bottomrule
\end{tabular}}
\begin{flushleft}\footnotesize\textit{Notes:} Monthly VAR in the unemployment gap and the market index, January 1973 to September 2026, 2020 months with dummies; unemployment first. $\Phi(h)$: response of the unemployment gap (percentage points) $h$ months after an innovation that raises the index by one unit. Granger: test that the lags of the index do not enter the unemployment equation. Interval: 5th and 95th percentiles of 500 block-bootstrap draws (12-month blocks of residuals). Months significant: first and last month with the 90 percent band above zero.\end{flushleft}
\end{table}
